
\renewcommand{\thechapter}{3}

\chapter{Quantitative Evaluation of Property-Based\\Testing}\label{chap:etna}

The second chapter of this thesis remarks a fragmentation caused by the design of
property-based testing libraries, leading to diverging features, separated ecosystems.
In this chapter, I focus on yet another fragmentation, this time caused by something
much more banal, a lack of general purpose evaluation benchmarks for PBT. The
literature is full of interesting novel techniques for enhancing random generation and
counterexample minimization, but these techniques are rarely evaluated within a
coherent context. Many times, the studies have relied on bespoke benchmarks
(\todo{examples}), if one were to build a web of studies and the benchmarks they used,
the result is somewhat reminiscent of the infamous public transport coverage of Ankara,
a small number of hub workloads that many studies reuse, many one-off benchmarks that
haven't been validated or reused in other studies, and a complete lack of
standardization.

This lack of common comparative analysis surely contributes to a lack of understanding
in the effects of these new techniques and changes; but the problem is persistent
beyond the literature and into the libraries themselves. All the property-based testing
libraries I have read in the course of developing this thesis are full of bespoke
decisions accummulated in the last 25 years. The choice of which sizing function to
use, the shrinking algorithm, the default settings on budgets and parameters... Many of
these decisions have never been subject to quantitative judgement beyond the initial
discussions, never been evaluated at scale, or understood in the broader context of PBT
usage.

In the beginning of my Ph.D. we set out to remediate this fragmentation and help build
a common quantitative understanding for the field. With Jessica and Harry implementing
the common evaluation framework in Python as well as the Haskell workloads, and I
working on the implementation of the Rocq pipeline and workloads, we have published
``Etna: An Evaluation Platform for Property-Based Testing (Experience
Report)''~\cite{etna-icfp23} in ICFP23.

After the initial publication, I have taken the lead role in future development of the
ETNA platform and horizontally grown it for OCaml with Nikhil Kamath, Racket with Ceren
Mert and Justine Frank, for Haskell with George Miao, and for Rust by myself. The
initial implementation as a library has been replaced with a standalone tool that could
be easily run cross-platform on many machines, eventually to the point we can start to
run our experiments as Github Actions. Many of these changes have found their ways into
publications, our JFP extension for the ICFP23 paper has incorporated the work by Ceren
and Nikhil as well as Tin from UPenn working together with Jessica.

The rest of this chapter includes excerpts edited from these papers, mainly the 2026
JFP extension of ETNA, with a major section on evaluation of shrinking from our
under-review work for the Haskell Symposium with George and Leo.

\section{ETNA}

ETNA is, as you may have thought of, named after the famous Sicilian volcano~\cn. The
name follows a convention established by prior work in the fuzz testing community,
Lava~\cite{Lava} and Magma~\cite{Hazimeh:2020:Magma}, both of which have been
instrumental to our design, particularly with the ground truth approach of Magma.

During the design and development of ETNA, we focused on a few key design principles,
centered around usefulness, extensibility, and maintainability.

\subsection{Evaluate for the ground truth, not for proxy metrics.}

How should an evaluation platform measure the effectiveness of a generator? The
software testing literature offers two common answers: \emph{code coverage} and
\emph{mutation testing}. Code coverage is attractive because it is easy to instrument
and compare, but it is at best an indirect signal: higher coverage does not always
translate to better bug finding~\cite{GopinathJG14,EvaluatingFuzzTesting}. For ETNA, we
instead chose mutation testing~\cite{JiaH11}, which measures the effectiveness of a
testing technique by artificially injecting mutations into the system under test and
checking whether testing is able to detect them.

Mutations used in prior evaluation work fall on a spectrum from manually sourced to
automatically synthesized~\cite{Hazimeh:2020:Magma,EvaluatingFuzzTesting,
    FixReverter,TestingNIjfp}. ETNA opts for manual sourcing. This choice makes the
benchmark harder to populate, since every task requires a meaningful mutant and a
property that exposes it, but it also lets us maintain a clearer ground truth: each
mutant should violate some aspect of the property specification. The workloads,
properties, and injected bugs used in this chapter are detailed in
\S~\ref{etna:workloads}.

\subsection{Use minimal, but precise interfaces.}

A key challenge during ETNA's development was that the platform needed to gather
comparable data about the testing effectiveness of a wide variety of existing PBT
frameworks, each of which reports such data in an ad-hoc and non-standardized manner.
Most frameworks report the number of inputs generated before a counterexample is found.
Far fewer expose timing statistics, and none of the frameworks we studied initially
offered a detailed breakdown of generation, testing, and minimization time. Similarly,
most frameworks report the number of inputs that fail to satisfy a precondition as
discards, while others, such as Rust's QuickCheck and Racket's RackCheck, do not
separate discards from passing tests.

These differences matter because ETNA is not evaluating a single library in isolation;
it is trying to compare workloads, strategies, frameworks, and languages. To make those
comparisons meaningful, we settled on a small set of metrics that are easy to measure
across frameworks and on a precise output format that implementations can validate
themselves: JSON conforming to a schema exposed by the ETNA
protocol.\footnote{\url{https://github.com/alpaylan/etna-cli/blob/main/PROTOCOL.md}}
The resulting interface is intentionally modest. It does not try to encode every
feature of every PBT library. Instead, it records the information that ETNA needs to
run experiments reproducibly and compare their outcomes.

\subsection{Every ETNA capability should be available to use manually.}

ETNA acts as an orchestration mechanism: it invokes testing tools, toggles mutations,
parses results, and performs analysis over the collected data. Such orchestration is
fragile even when every component is individually well engineered, because ETNA sits
between loosely coupled systems that were developed independently. In practice, any
opaqueness in this process can make failures difficult to diagnose and, in the worst
case, unrecoverable.

The final guiding principle of ETNA is therefore transparency. Each step of an
experiment should correspond to an operation that a user can understand and, when
needed, reproduce manually. This principle shaped the transition from the initial
library implementation to the standalone command-line tool described in the next
section. Rather than hiding build steps, mutation activation, test execution, and
analysis behind a single opaque pipeline, ETNA exposes them as ordinary commands. That
design makes the platform easier to debug, easier to extend to new languages, and
easier to trust when an experiment produces surprising results.

\section{ETNA Tools}

\section{Workloads} \label{etna:workloads}
\section{Experiments}
\subsection{Evaluating Random Generation}
\subsubsection{Rocq}
\subsubsection{OCaml}
\subsection{Evaluating Shrinking}
\section{Discussion}
\section{Future Work}
